\section{Introduction}
\label{sec:introduction}

Single-inclusive hadron production probes short-distance QCD together with
parton distributions and fragmentation functions.  At fixed order, its hard
parts contain several structures that cannot be reconstructed from values at
isolated phase-space points: Laurent poles in dimensional regularization,
noninteger powers at the real-emission endpoint, distributions in the
partonic variable $w$, and, for polarized observables, evanescent BMHV
tensors.  An analytic calculation must preserve these structures from the
amplitude to the final coefficient.

The purpose of this work is to organize that calculation as one connected
chain.  Feynman amplitudes and their conjugates are projected onto the
leading-twist collinear density matrices.  The real phase space is represented
by positive-energy cuts, and literal propagators are partial fractioned before
integral families are chosen.  Cut-aware integration-by-parts (IBP) identities
reduce the result to master integrals.  Their coefficients are reconstructed
exactly over finite fields, while the masters themselves are evaluated as
analytic functions by direct parametric integration or by differential
equations with analytic boundary data.  Numerical evaluations enter only as
independent comparisons.

The NLO real-emission problem provides a complete analytic example of the
master-integral part of this chain.  Its six masters reduce to one cut bubble
and one generic top function evaluated at five kinematic arguments.  A
triangular differential system fixes that function from one boundary value,
and its connection formula separates the endpoint branches before the
$\eps$ expansion.  For the NNLO double-real sector considered here, exact
IBP reduction gives 342 masters.  We distinguish powered-integral identities,
denominator identities, differential coupling, and equality of boundary
periods; these notions answer different questions and lead to different
counts.  One genuinely coupled eight-master family is solved through
$\mathcal O(\eps^2)$ as a worked NNLO example.

The paper follows the calculation rather than the software layout.
Section~\ref{sec:factorized-observable} defines the observable, kinematics,
polarization projections, and hard coefficient.
Section~\ref{sec:cut-families} derives the cut families from ordered amplitude
interferences.  Section~\ref{sec:ibp-coefficients} discusses IBP reduction and
exact rational master coefficients.  Section~\ref{sec:analytic-masters}
describes the analytic evaluation of the cut masters.
Sections~\ref{sec:nlo-example} and \ref{sec:nnlo-double-real} give the NLO and
NNLO examples.  Reproducibility conditions and the present scope are stated in
Section~\ref{sec:reproducibility}.
